\begin{table}[h!]
\caption{Dynamic Olley-Pakes decomposition for ARPK growth, \% , 2001-2008}
\label{table:data_op_decomp}
    \centering
\scalebox{0.7}{\begin{tabular}{l*{7}{c}}
\hline\hline
&  Overall &   Incumbent &  Distributional &  Reallocation  &  Newborn  &   Exiting  \\
\hline
Full sample &        14 &        17 &         3 &        15 &        -3 &         0 \\
Firms with foreign owners &        17 &        17 &        11 &         5 &        -2 &         2 \\
Firms without foreign owners &        12 &        19 &         1 &        18 &        -3 &        -4 \\
Exporters with foreign owners &        21 &        14 &         7 &         7 &        -1 &         8 \\
Exporters without foreign owners &         3 &         4 &       -16 &        20 &        12 &       -12 \\
Non-exporters with foreign owners &       -11 &        35 &        20 &        15 &       -26 &       -20 \\
Non-exporters without foreign owners &        20 &        41 &         7 &        33 &        -4 &       -16 \\
\hline
\hline
\end{tabular}}
\caption*{  \parbox{0.95\textwidth}{\scriptsize \textit{Note:}  Capital used as weights $s_j = \frac{k_j}{K}$ used to decompose the overall growth of average revenue productivity of capital, $ARPK = \sum_{j \in J} ARPK_j s_j$ with $ARPK_j = \frac{VA_j}{K_j}$. Overall refers to all firms while Incumbent, Newborn and Exiting correspond to separating the set of indices $J$ into firms that are either incumbents in, enter to, or exit from the sub-sample year on year. A firm can be incumbent in one year, exiting in another one. Incumbent ARPK gains can be due to firm-level distributional shifts or as a result of capital reallocation. I provide further details about the interpretation of each column in Appendix \ref{appendix:firm_data}}}
\end{table}
